% Part: lambda-calculus
% Chapter: introduction
% Section: syntax

\documentclass[../../../include/open-logic-section]{subfiles}

\begin{document}

\olfileid{lam}{int}{syn}
\olsection{લૅમ્બડા કલનની વાક્યરચના}

સૌપ્રથમ $x$, $y$, $z$,~\dots{} ચલોની શ્રેણી અને
$a$, $b$, $c$,~\dots{} જેવા કેટલાક અચલ સંકેતો લઈએ. પદોનો ગણ
નીચે પ્રમાણે આગમનથી વ્યાખ્યાયિત થાય છે:
\begin{enumerate}
\item દરેક ચલ પદ છે.
\item દરેક અચલ પદ છે.
\item $M$ અને $N$ પદો હોય તો $(MN)$ પણ પદ છે.
\item $M$ પદ અને $x$ ચલ હોય તો $(\lambd[x][M])$ પણ પદ છે.
\end{enumerate}
$(MN)$ સ્વરૂપનાં પદોને \emph{અનુપ્રયોગો} અને
$(\lambd[x][M])$ સ્વરૂપનાં પદોને \emph{અમૂર્તીકરણો} કહે છે.

કોઈપણ અચલ વિનાની પદ્ધતિને \emph{શુદ્ધ} લૅમ્બડા કલન કહે છે.
મુખ્યત્વે આપણે શુદ્ધ~$\lambd$-કલનમાં કામ કરીશું; તેથી બધાં નાના
અક્ષરો ચલો દર્શાવશે. મોટા અક્ષરો ($M$, $N$ વગેરે) $\lambd$-કલનનાં
પદો દર્શાવશે.

આપણે સંકેતન અંગેના કેટલાક નિયમો અનુસરીશું:

\begin{conv}
\begin{enumerate}
\item કૌંસ છોડી દેવાય ત્યારે અનુપ્રયોગ ડાબેથી જમણે થાય છે.
  દાખલા તરીકે, $M$, $N$, $P$ અને $Q$ પદો હોય તો~$MNPQ$ એ
  $(((MN)P)Q)$નું સંક્ષિપ્ત સ્વરૂપ છે.
\item ફરી, કૌંસ છોડી દેવાય ત્યારે લૅમ્બડા અમૂર્તીકરણનો વ્યાપ
  શક્ય તેટલો મોટો લેવાય છે. દાખલા તરીકે, $\lambd[x][MNP]$ને
  $(\lambd[x][((MN)P)])$ એમ વાંચવું.
\item એક લૅમ્બડાથી અનેક ચલોનું અમૂર્તીકરણ કરી શકાય છે.
  દાખલા તરીકે, $\lambd[xyz][M]$ એ
  $\lambd[x][\lambd[y][\lambd[z][M]]]$નું સંક્ષિપ્ત સ્વરૂપ છે.
\end{enumerate}
\end{conv}

દાખલા તરીકે,
\[
\lambd[xy][xxyx \lambd[z][xz]]
\]
એ નીચેનું સંક્ષિપ્ત સ્વરૂપ છે:
\[
\lambd[x][\lambd[y][((((xx)y)x)(\lambd[z][(xz)]))]].
\]
આ નિયમો યાદ રાખવા જોઈએ. શરૂઆતમાં એ માથું ખાઈ જશે, પરંતુ
ધીમે ધીમે ટેવ પડી જશે; પછી અનેક કૌંસોની ગૂંચવણ સંભાળવા કરતાં
આ નિયમો ઓછા કંટાળાજનક લાગશે.

માત્ર બંધ ચલોની નામાવલિમાં જુદાં પડતાં બે પદોને
$\alpha$-સામ્યવાળાં કહે છે; દાખલા તરીકે, $\lambd[x][x]$ અને
$\lambd[y][y]$. તેમને ``એક જ'' પદ માનવું અનુકૂળ રહેશે.
એટલે આપણે~$M$ અને~$N$ એક જ છે એમ કહીએ ત્યારે બંધ ચલોનું નામ
ફેરવવા સુધીની સમાનતા પણ તેમાં સામેલ છે. કોઈ~$\lambd$ના વ્યાપમાં
આવતા ચલોને ``બંધ'' અને બાકીનાને ``મુક્ત'' કહે છે. અગાઉના ઉદાહરણમાં
મુક્ત ચલો નથી; પરંતુ
\[
(\lambd[z][yz])x
\]
માં~$y$ અને~$x$ મુક્ત છે, જ્યારે~$z$ બંધ છે.

\end{document}
